% TOP_PROBLEM: {"id":30007544,"problem_number":"LOCAL-30007544","title":"AI and aggregate productivity","category_id":21,"category":{"id":21,"name":"economics","display_name":"Economics","description":"Open economic theory statements, published research agendas, and proposed scoped specifications; question provenance and historical priority are qualified per record.","slug":"economics","order_index":21,"created_at":"2026-10-08T00:00:00Z"},"status":"research_question","external_url":null,"legacy_ids":[],"published":false,"statement":"Identify an aggregate productivity effect of a precisely defined AI capability and adoption path, incorporating complementary inputs and equilibrium adjustments.","economics":{"original_id":33,"field":"Growth and productivity","kind":"identification","formalization_basis":"adapted_research_question","source_status":"research_agenda","priority_status":"earliest_found","priority_note":"This is the earliest explicit relevant statement verified in this audit, not a claim of original priority; older literature and earlier versions may contain antecedents.","earliest_verified_reference":{"authors":["Erik Brynjolfsson","Anton Korinek","Ajay Agrawal"],"title":"The Economics of Transformative AI: A Research Agenda","year":2024,"url":"https://digitaleconomy.stanford.edu/app/uploads/2024/11/ETAI-White-Paper.pdf","doi":null,"locator":"Section 3.1, “Economic Growth,” printed pp. 8–9; 23 October 2024 draft","evidence_summary":"The agenda explicitly asks how transformative AI affects growth, how adoption and bottlenecks shape its impact, and what economic mechanisms constrain the transition. The quantitative treatment definition below is newly specified."},"current_reference":{"authors":["Erik Brynjolfsson","Anton Korinek","Ajay Agrawal"],"title":"The Economics of Transformative AI: A Research Agenda","year":2024,"url":"https://digitaleconomy.stanford.edu/app/uploads/2024/11/ETAI-White-Paper.pdf","doi":null,"locator":"Section 3.1, “Economic Growth,” printed pp. 8–9; 23 October 2024 draft","evidence_summary":"The agenda explicitly asks how transformative AI affects growth, how adoption and bottlenecks shape its impact, and what economic mechanisms constrain the transition. The quantitative treatment definition below is newly specified."},"formalization_note":"The source explicitly sets an AI growth research agenda. The capability-defined potential-outcome estimand is an adaptation, not a source-stated numerical conjecture."},"statusQualification":"Published question or research agenda verified at the cited locator. The mathematical specification is adapted for this catalogue and includes additional assumptions; source publication does not establish first-ever priority or certify this exact model remains unsolved. The source explicitly sets an AI growth research agenda. The capability-defined potential-outcome estimand is an adaptation, not a source-stated numerical conjecture.","statusReviewedAt":"2026-10-08"}
% FIRST_POSTED: 2026-10-08
% ENTRY_KIND: statement_only
% !TeX program = lualatex
\documentclass[11pt]{article}
\usepackage{fontspec}
\usepackage[a4paper,margin=25mm]{geometry}
\usepackage{amsmath,amssymb,mathtools}
\usepackage{xurl}
\usepackage[hidelinks]{hyperref}
\newcommand{\sref}[2]{\hyperlink{source:#1:#2}{[S#2]}}
\setlength{\emergencystretch}{3em}
\begin{document}
% problemId:30007544
\section*{AI and aggregate productivity}\label{30007544}
\subsection{Definitions and mathematical statement}
Fix an economy with firms \(i\), occupations \(o\), and tasks \(j\). Let \(a_{ijt}\) denote use of a precisely defined AI capability, \(K^C_{it}\) complementary organizational capital, and \(K^A_{it}\) compute capital. Production combines task outputs through a declared aggregation technology; task demand, wages, entry, and investment respond in general equilibrium. Specify an AI-feasible counterfactual path \(a^1\) and a restricted path \(a^0\) in which that capability is unavailable, holding other primitive innovations fixed.
For measured aggregate total factor productivity \(A_t\) and horizon \(H\), define
\[
\Delta_H^{AI}=E[\log A_{t+H}(a^1)-\log A_{t+H}(a^0)].
\]
Identify a credible set for \(\Delta_H^{AI}\) using task experiments, adoption instruments, and independently measured complementary investment, with explicit assumptions linking microeconomic responses to equilibrium aggregation. Separate direct task efficiency from adoption selection, displaced labor, new task creation, and measurement of intangible capital. Impose observed constraints on compute, energy, managerial capacity, and diffusion rather than extrapolating laboratory task improvements mechanically. Validate transition predictions against subsequent adoption cohorts. The target concerns a declared capability, counterfactual, and horizon; forecasts about unspecified future intelligence or unconstrained autonomous research are not identified by current firm-level experiments.

\par\textbf{Formulation and scope.} The source explicitly sets an AI growth research agenda. The capability-defined potential-outcome estimand is an adaptation, not a source-stated numerical conjecture.

\par\textbf{Open-question provenance.} An explicit question or research agenda was verified at the source locator. This does not by itself certify that every later version remains unresolved \sref{30007544}{1}.

\par\textbf{Historical priority.} This is the earliest explicit relevant statement verified in this audit, not a claim of original priority; older literature and earlier versions may contain antecedents.

\subsection{Short English statement}
Identify an aggregate productivity effect of a precisely defined AI capability and adoption path, incorporating complementary inputs and equilibrium adjustments.

\subsection{Sources}
\begin{itemize}\raggedright
\item[\textnormal{[S1]}]\hypertarget{source:30007544:1}{} Erik Brynjolfsson, Anton Korinek, Ajay Agrawal. 2024. The Economics of Transformative AI: A Research Agenda. Section 3.1, “Economic Growth,” printed pp. 8–9; 23 October 2024 draft.
\url{https://digitaleconomy.stanford.edu/app/uploads/2024/11/ETAI-White-Paper.pdf}.
{\small\textit{Earliest verified question/agenda reference; historical priority qualified below. The agenda explicitly asks how transformative AI affects growth, how adoption and bottlenecks shape its impact, and what economic mechanisms constrain the transition. The quantitative treatment definition below is newly specified.}}
\end{itemize}
\end{document}
